\documentclass[11pt,a4paper]{scrartcl}

\usepackage{amssymb}
\usepackage{amsmath}

\usepackage[utf8]{inputenc}
\usepackage[T1]{fontenc}
\newcommand{\ats}[1]{\par\noindent\textit{Atsakymas:} #1}
\usepackage{cancel}
\usepackage{tikz}
\usepackage{lmodern} 
\usepackage{amsmath,amssymb,amsthm}
\usepackage{graphicx}
\usepackage[dvipsnames,svgnames]{xcolor}
\usepackage{hyperref}
\usepackage{mathrsfs}
\usepackage[shortlabels]{enumitem}
\usepackage[obeyFinal,textsize=scriptsize,shadow,loadshadowlibrary]{todonotes}
\usepackage{mathtools}
\usepackage{microtype}

%%% Paragraph layout
\setlength{\parindent}{0pt}
\setlength{\parskip}{0.6\baselineskip}

%%% Colored sections
\renewcommand*{\sectionformat}%
  {\color{purple}\S\thesection\enskip}
\renewcommand*{\subsectionformat}%
  {\color{purple}\S\thesubsection\enskip}
\renewcommand*{\subsubsectionformat}%
  {\color{purple}\S\thesubsubsection\enskip}
\addtokomafont{paragraph}{\color{orange!35!black}\P\ }
\KOMAoptions{numbers=noenddot}

%%% Title
\addtokomafont{subtitle}{\Large}
\setkomafont{author}{\Large\scshape}
\setkomafont{date}{\Large\normalsize}
% Neigiamas vertikalus tarpas, kuris pakelia dokumento antraštę į viršų. 
% Galite keisti -2.5cm į -3cm (aukščiau) arba -1.5cm (žemiau).
\titlehead{\vspace*{-7.5cm}} 

% PRIDĖTA: Automatiškai perrašome datą į mokyklos pavadinimą
\AtBeginDocument{%
  \date{\normalsize Vilniaus licėjus}
}

%%% Page setup
\usepackage[headsepline]{scrlayer-scrpage}
\renewcommand{\headfont}{}
\addtolength{\textheight}{3.14cm}
\setlength{\footskip}{0.5in}
\setlength{\headsep}{10pt}
% Puslapio antraštė turi rodyti tik konkrečius dokumento duomenis.
% Nustatome juos aiškiai, kad neatsirastų "\@author" / "\@title" arba senų reikšmių.
\makeatletter
\AtBeginDocument{%
  \ihead{\footnotesize\textbf{Andrius Berniukevičius}}%
  \ohead{\footnotesize\textbf{\@title}}%
}
\makeatother
\automark{section}
\chead{}
\cfoot{\pagemark}

%%% Macros
\providecommand{\ol}{\overline}
\providecommand{\ul}{\underline}
\providecommand{\wt}{\widetilde}
\providecommand{\wh}{\widehat}
\providecommand{\eps}{\varepsilon}
\providecommand{\half}{\frac{1}{2}}
\providecommand{\inv}{^{-1}}
\newcommand{\dang}{\measuredangle} %% Directed angle
\newcommand{\vocab}[1]{\textbf{\color{blue}\sffamily #1}}
\providecommand{\CC}{\mathbb C}
\providecommand{\FF}{\mathbb F}
\providecommand{\NN}{\mathbb N}
\providecommand{\QQ}{\mathbb Q}
\providecommand{\RR}{\mathbb R}
\providecommand{\ZZ}{\mathbb Z}
\providecommand{\ts}{\textsuperscript}
\providecommand{\dg}{^\circ}
\providecommand{\ii}{\item}
\providecommand{\defeq}{\coloneqq}
\DeclareMathOperator*{\lcm}{lcm}
\DeclareMathOperator*{\argmin}{arg min}
\DeclareMathOperator*{\argmax}{arg max}
\providecommand{\hrulebar}{\par
\hspace{\fill}\rule{0.95\linewidth}{.7pt}\hspace{\fill}
\par\nointerlineskip \vspace{\baselineskip}}

%%% Optional colored theorems
\usepackage{thmtools}
\usepackage[framemethod=TikZ]{mdframed}
\mdfdefinestyle{mdbluebox}{%
  roundcorner=10pt,
  linewidth=1pt,
  skipabove=12pt,
  innerbottommargin=9pt,
  skipbelow=2pt,
  linecolor=blue,
  nobreak=true,
  backgroundcolor=TealBlue!5,
}
\declaretheoremstyle[
  headfont=\sffamily\bfseries\color{MidnightBlue},
  mdframed={style=mdbluebox},
  headpunct={\\[3pt]},
  postheadspace={0pt}
]{thmbluebox}
\mdfdefinestyle{mdredbox}{%
  linewidth=0.5pt,
  skipabove=12pt,
  frametitleaboveskip=5pt,
  frametitlebelowskip=0pt,
  skipbelow=2pt,
  frametitlefont=\bfseries,
  innertopmargin=4pt,
  innerbottommargin=8pt,
  nobreak=true,
  backgroundcolor=Salmon!5,
  linecolor=RawSienna,
}
\declaretheoremstyle[
  headfont=\bfseries\color{RawSienna},
  mdframed={style=mdredbox},
  headpunct={\\[3pt]},
  postheadspace={0pt},
]{thmredbox}
\mdfdefinestyle{mdgreenbox}{%
  skipabove=8pt,
  linewidth=2pt,
  rightline=false,
  leftline=true,
  topline=false,
  bottomline=false,
  linecolor=ForestGreen,
  backgroundcolor=ForestGreen!5,
}
\declaretheoremstyle[
  headfont=\bfseries\sffamily\color{ForestGreen!70!black},
  bodyfont=\normalfont,
  spaceabove=2pt,
  spacebelow=1pt,
  mdframed={style=mdgreenbox},
  headpunct={ --- },
]{thmgreenbox}
\mdfdefinestyle{mdblackbox}{%
  skipabove=8pt,
  linewidth=3pt,
  rightline=false,
  leftline=true,
  topline=false,
  bottomline=false,
  linecolor=black,
  backgroundcolor=RedViolet!5!gray!5,
}
\declaretheoremstyle[
  headfont=\bfseries,
  bodyfont=\normalfont\small,
  spaceabove=0pt,
  spacebelow=0pt,
  mdframed={style=mdblackbox}
]{thmblackbox}

%% Aplinkos su rėmeliais (thmtools)
\declaretheorem[style=thmbluebox,name=Teorema]{theorem}
\declaretheorem[style=thmbluebox,name=Lema]{lemma}
\declaretheorem[style=thmbluebox,name=Savybė]{property}
\declaretheorem[style=thmbluebox,name=Teiginys]{proposition}
\declaretheorem[style=thmbluebox,name=Išvada]{corollary}
\declaretheorem[style=thmbluebox,name=Teorema,numbered=no]{theorem*}
\declaretheorem[style=thmbluebox,name=Lema,numbered=no]{lemma*}
\declaretheorem[style=thmbluebox,name=Teiginys,numbered=no]{proposition*}
\declaretheorem[style=thmbluebox,name=Išvada,numbered=no]{corollary*}
\declaretheorem[style=thmgreenbox,name=Algoritmas]{algorithm}
\declaretheorem[style=thmgreenbox,name=Algoritmas,numbered=no]{algorithm*}
\declaretheorem[style=thmgreenbox,name=Teiginys]{claim}
\declaretheorem[style=thmgreenbox,name=Teiginys,numbered=no]{claim*}
\declaretheorem[style=thmredbox,name=Pavyzdys]{example}
\declaretheorem[style=thmredbox,name=Pavyzdys,numbered=no]{example*}
\declaretheorem[style=thmblackbox,name=Pastaba]{remark}
\declaretheorem[style=thmblackbox,name=Pastaba,numbered=no]{remark*}

%% Standartinės amsthm aplinkos (Apibrėžimai ir užduotys)
\theoremstyle{definition}
\ifcsname conjecture\endcsname\else\newtheorem{conjecture}{Hipotezė}\fi
\ifcsname definition\endcsname\else\newtheorem{definition}{Apibrėžimas}\fi
\ifcsname fact\endcsname\else\newtheorem{fact}{Faktas}\fi
\ifcsname answer\endcsname\else\newtheorem{answer}{Atsakymas}\fi
\ifcsname ques\endcsname\else\newtheorem{ques}{Klausimas}\fi
\ifcsname exercise\endcsname\else\newtheorem{exercise}{Užduotis}\fi
\ifcsname problem\endcsname\else\newtheorem{problem}{Uždavinys}[section]\fi
\ifcsname conjecture*\endcsname\else\newtheorem*{conjecture*}{Hipotezė}\fi
\ifcsname definition*\endcsname\else\newtheorem*{definition*}{Apibrėžimas}\fi
\ifcsname fact*\endcsname\else\newtheorem*{fact*}{Faktas}\fi
\ifcsname answer*\endcsname\else\newtheorem*{answer*}{Atsakymas}\fi
\ifcsname ques*\endcsname\else\newtheorem*{ques*}{Klausimas}\fi
\ifcsname exercise*\endcsname\else\newtheorem*{exercise*}{Užduotis}\fi
\ifcsname problem*\endcsname\else\newtheorem*{problem*}{Uždavinys}\fi

\newcounter{answerssection}
\newcommand{\answerssection}{%
  \setcounter{answerssection}{\value{section}}%
  \section{Atsakymai}%
  \setcounter{problem}{0}%
  \renewcommand{\theproblem}{\arabic{answerssection}.\arabic{problem}}%
}

%% GALUTINIS SPRENDIMAS PROOF APLINKAI
%% --- PRADŽIA: Įrodymo aplinkos sutvarkymas ---
\makeatletter

% 1. Užtikriname, kad žodis būtų "Įrodymas", net jei "babel" paketas bando jį perrašyti
\AtBeginDocument{%
  \renewcommand{\proofname}{Įrodymas}%
  \@ifpackageloaded{babel}{%
    \addto\captionslithuanian{\renewcommand{\proofname}{Įrodymas}}%
  }{}%
}

% 2. Priverstinai pašaliname scrartcl klasės bold šriftą ir paliekame TIK italic
% 3. Sujungiame su tuščiu kvadratėliu dešiniajame kampe.
\renewcommand{\qedsymbol}{\ensuremath{\Box}}
\renewenvironment{proof}[1][\proofname]{\par
  \pushQED{\qed}%
  \normalfont \topsep6\p@\@plus6\p@\relax
  \trivlist
  \item[\hskip\labelsep
        \normalfont\itshape % Priverčia naudoti tik pasvirą šriftą, kaip nuotraukoje
    #1\@addpunct{.}]\ignorespaces
}{%
  \popQED\endtrivlist\@endpefalse
}
\newenvironment{solution}{\par
  \normalfont \topsep6\p@\@plus6\p@\relax
  \trivlist
  \item[\hskip\labelsep
        \normalfont\itshape
    \textit{Sprendimas}\@addpunct{.}]\ignorespaces
}{%
  \endtrivlist\@endpefalse
}
\newcommand{\ans}[1]{\par\noindent\textit{Atsakymas:} #1}
\makeatother


\title{Algebra I: Lygčių sistemos su dviem kintamaisiais}
\author{Andrius Berniukevičius}

\begin{document}

\maketitle

\section{Keitimo metodo krachas}

Mokyklinėse lygčių sistemose dažnai pakanka vieną kintamąjį išreikšti per kitą ir jį įstatyti į likusią lygtį. Olimpiadinėse sistemose toks metodas ne visuomet yra veiksmingas: jis gali sukurti aukštų laipsnių lygtis arba nepatogias šaknines išraiškas.

Įsivaizduokite sistemą:
\[
\left\{\begin{aligned}
    x^2 + y^2 &= 13 \\
    x^3 + y^3 &= 35
\end{aligned}\right.
\]
Jei iš pirmosios lygties išreikštume $y = \sqrt{13-x^2}$ ir šią išraišką įstatytume į antrąją lygtį, gautume gerokai sudėtingesnę lygtį su šaknimis ir aukštais laipsniais.

Todėl svarbu atpažinti sistemos struktūrą ir pasirinkti jai tinkamą metodą: simetrinius reiškinius, pilnųjų kvadratų išskyrimą arba skaidymą dauginamaisiais.

% ==========================================
% BAZINIAI GINKLAI
% ==========================================
\section{Pagrindiniai metodai}

\subsection{Simetriniai reiškiniai}
Jei sukeitus $x$ ir $y$ vietomis sistemos lygtys nepasikeičia, sistema yra \vocab{simetrinė}. Tuomet patogu ją perrašyti per kintamųjų sumą ir sandaugą:
$$ u = x + y $$
$$ v = x y $$
Naudojame formules:
$$ x^2 + y^2 = (x+y)^2 - 2xy = u^2 - 2v $$
$$ x^3 + y^3 = (x+y)(x^2 - xy + y^2) = u(u^2 - 3v) $$

\begin{example}[Simetrinė sistema]
Išspręskite lygčių sistemą:
\[
\left\{\begin{aligned}
    x + y + xy &= 11 \\
    x^2 y + x y^2 &= 30
\end{aligned}\right.
\]
\end{example}

\begin{solution}
Sukeitus $x$ ir $y$ vietomis, abi sistemos lygtys nepasikeičia. Todėl patogu naudoti simetrinius reiškinius
\[ u=x+y, \qquad v=xy. \]
Pirmoji lygtis tampa $u+v=11$. Antrojoje iškeliame $xy$:
\[ x^2y+xy^2=xy(x+y)=uv. \]
Taigi pradinę sistemą perrašome taip:
\[
\left\{\begin{aligned}
    u+v&=11 \\
    uv&=30
\end{aligned}\right.
\]

Ieškome dviejų skaičių, kurių suma yra $11$, o sandauga $30$. Tai yra skaičiai $5$ ir $6$, todėl galimi du atvejai:
\[
\text{1) } u=5,\ v=6, \qquad \text{2) } u=6,\ v=5.
\]

Pirmuoju atveju turime $x+y=5$ ir $xy=6$. Pagal Vijeto teoremą $x$ ir $y$ yra lygties
\[
t^2-5t+6=0
\]
šaknys. Išskaidę gauname $(t-2)(t-3)=0$, todėl $(x,y)=(2,3)$ arba $(3,2)$.

Antruoju atveju turime $x+y=6$ ir $xy=5$. Dabar $x$ ir $y$ yra lygties
\[
t^2-6t+5=0
\]
šaknys. Išskaidę gauname $(t-1)(t-5)=0$, todėl $(x,y)=(1,5)$ arba $(5,1)$.

Visos gautos poros tenkina pradinę sistemą, todėl kitų realių sprendinių nėra.

\ats{$(2,3), (3,2), (1,5), (5,1)$}
\end{solution}

\subsection{Pilno kvadrato išskyrimas}
Kartais dvi atskiros lygtys atrodo negražiai, bet jas \textbf{sudėjus} arba \textbf{atėmus}, iš niekur nieko atsiranda tobulas pilnas kvadratas $(x \pm y)^2$. Jūsų akys turi išmokti ieškoti kvadratų ir dvigubų sandaugų ($2xy$) porų.

\begin{example}[Pilno kvadrato konstravimas]
Išspręskite lygčių sistemą:
\[
\left\{\begin{aligned}
    x^2 + y^2 &= 25 \\
    xy &= 12
\end{aligned}\right.
\]
\end{example}

\begin{solution}
Šioje sistemoje turime kvadratų sumą $x^2+y^2$ ir sandaugą $xy$. Padauginę antrąją lygtį iš $2$, gauname
\[
2xy=24.
\]

Sudedame šią lygtį su pirmąja:
\[
x^2+y^2+2xy=25+24,
\]
\[
(x+y)^2=49.
\]
Todėl
\[
x+y=7 \quad\text{arba}\quad x+y=-7.
\]

Iš pirmosios lygties atimame $2xy=24$:
\[
x^2+y^2-2xy=25-24,
\]
\[
(x-y)^2=1.
\]
Todėl
\[
x-y=1 \quad\text{arba}\quad x-y=-1.
\]

Išsprendžiame keturias galimas tiesines sistemas:
\begin{align*}
    x+y=7,\ x-y=1 &\implies (x,y)=(4,3), \\
    x+y=7,\ x-y=-1 &\implies (x,y)=(3,4), \\
    x+y=-7,\ x-y=1 &\implies (x,y)=(-3,-4), \\
    x+y=-7,\ x-y=-1 &\implies (x,y)=(-4,-3).
\end{align*}

Kiekviena pora tenkina abi pradines lygtis, todėl gauname visus sprendinius.

\ats{$(4,3), (3,4), (-4,-3), (-3,-4)$}
\end{solution}

\subsection{Skaidymas dauginamaisiais}
Jei vienoje lygties pusėje matote nulį (pvz., $x^2 - xy - 2y^2 = 0$), tai yra galingas signalas, kad kairę pusę reikia skaidyti dauginamaisiais. Gavę sandaugą, lygią nuliui, sistemą padalinsite į du labai paprastus atskirus atvejus.

\begin{example}[Skaidymas dauginamaisiais]
Išspręskite lygčių sistemą:
\[
\left\{\begin{aligned}
    x^2 - xy - 2y^2 &= 0 \\
    x^2 + 2y^2 &= 18
\end{aligned}\right.
\]
\end{example}

\begin{solution}
Pirmoji lygtis lygi nuliui, todėl jos kairiąją pusę skaidome dauginamaisiais:
\begin{align*}
x^2-xy-2y^2 &= x^2+xy-2xy-2y^2 \\
             &= x(x+y)-2y(x+y) \\
             &= (x+y)(x-2y).
\end{align*}
Taigi gauname:
\[
(x+y)(x-2y)=0.
\]

Sandauga lygi nuliui tada, kai bent vienas jos dauginamasis lygus nuliui. Todėl nagrinėjame du atvejus.

\textbf{1 atvejis:} $x+y=0$, taigi $x=-y$. Įstatome į antrąją pradinę lygtį:
\[
(-y)^2+2y^2=18 \implies 3y^2=18 \implies y^2=6.
\]
Vadinasi, $y=\pm\sqrt6$. Atitinkamai gauname:
\[
(x,y)=(-\sqrt6,\sqrt6) \quad\text{arba}\quad (\sqrt6,-\sqrt6).
\]

\textbf{2 atvejis:} $x-2y=0$, taigi $x=2y$. Įstatome į antrąją lygtį:
\[
(2y)^2+2y^2=18 \implies 6y^2=18 \implies y^2=3.
\]
Vadinasi, $y=\pm\sqrt3$. Atitinkamai gauname:
\[
(x,y)=(2\sqrt3,\sqrt3) \quad\text{arba}\quad (-2\sqrt3,-\sqrt3).
\]

Šie keturi sprendiniai apima abu pirmosios lygties išskaidymo atvejus.

\ats{$(-\sqrt6,\sqrt6), (\sqrt6,-\sqrt6), (2\sqrt3,\sqrt3), (-2\sqrt3,-\sqrt3)$}
\end{solution}

\vspace{0.5cm}

% ==========================================
% UŽDAVINIAI SAVARANKIŠKAM SPRENDIMUI
% ==========================================
\newpage
\section{Uždaviniai savarankiškam sprendimui (20 iššūkių)}
Išspręskite lygčių sistemas:

\begin{enumerate}
    \item $\left\{\begin{aligned} x+y&=7 \\ xy&=12 \end{aligned}\right.$
    \item $\left\{\begin{aligned} x-y&=2 \\ xy&=15 \end{aligned}\right.$
    \item $\left\{\begin{aligned} x^2+y^2&=25 \\ x-y&=1 \end{aligned}\right.$
    \item $\left\{\begin{aligned} \frac{1}{x} + \frac{1}{y} &= \frac{5}{6} \\ x+y&=5 \end{aligned}\right.$

    \item $\left\{\begin{aligned} x^2+y^2&=34 \\ xy&=15 \end{aligned}\right.$
    \item $\left\{\begin{aligned} x^2+y^2&=58 \\ x+y&=10 \end{aligned}\right.$
    \item $\left\{\begin{aligned} x^2-y^2&=9 \\ x-y&=1 \end{aligned}\right.$
    \item $\left\{\begin{aligned} x^2y - xy^2&=6 \\ x-y&=1 \end{aligned}\right.$
    \item $\left\{\begin{aligned} x+y+xy&=5 \\ x^2+y^2&=5 \end{aligned}\right.$
    \item $\left\{\begin{aligned} x^3+y^3&=35 \\ x+y&=5 \end{aligned}\right.$
    \item $\left\{\begin{aligned} \frac{x}{y} + \frac{y}{x} &= \frac{13}{6} \\ x+y&=5 \end{aligned}\right.$
    \item $\left\{\begin{aligned} x^2+y^2+x+y&=18 \\ xy&=6 \end{aligned}\right.$
    \item $\left\{\begin{aligned} x^3-y^3&=19 \\ x-y&=1 \end{aligned}\right.$
    \item $\left\{\begin{aligned} x^2-4y^2&=0 \\ x^2+y^2&=5 \end{aligned}\right.$
    \item $\left\{\begin{aligned} x^2-3xy+2y^2&=0 \\ 2x^2-y^2&=7 \end{aligned}\right.$
    \item $\left\{\begin{aligned} 2x^2 - 3xy + y^2&=0 \\ x^2-y^2&=-3 \end{aligned}\right.$
    \item $\left\{\begin{aligned} x^4+y^4&=17 \\ x+y&=3 \end{aligned}\right.$
    \item $\left\{\begin{aligned} x^2y+xy^2&=20 \\ x^3+y^3&=65 \end{aligned}\right.$
    \item $\left\{\begin{aligned} (x+y)(x^2+y^2)&=15 \\ (x-y)(x^2-y^2)&=3 \end{aligned}\right.$
    \item Raskite sveikuosius lygčių sistemos sprendinius:$$\left\{\begin{aligned} x^2+y^2&=10 \\ x^3+y^3&=26 \end{aligned}\right.$$
\end{enumerate}

% ==========================================
% SAVARANKIŠKŲ UŽDAVINIŲ SPRENDIMAI
% ==========================================
% ==========================================
% ATSAKYMAI
% ==========================================
\newpage
\section*{Atsakymai}

\begin{enumerate}
    \item $(3,4), (4,3)$
    \item $(5,3), (-3,-5)$
    \item $(4,3), (-3,-4)$
    \item $(2,3), (3,2)$
    \item $(5,3), (3,5), (-5,-3), (-3,-5)$
    \item $(3,7), (7,3)$
    \item $(5,4)$
    \item $(3,2), (-2,-3)$
    \item $(1,2), (2,1)$
    \item $(2,3), (3,2)$
    \item $(2,3), (3,2)$
    \item $(2,3), (3,2), (-3+\sqrt{3}, -3-\sqrt{3}), (-3-\sqrt{3}, -3+\sqrt{3})$
    \item $(3,2), (-2,-3)$
    \item $(2,1), (-2,-1), (-2,1), (2,-1)$
    \item $(\sqrt{7/3},\sqrt{7/3}), (-\sqrt{7/3},-\sqrt{7/3}), (2,1), (-2,-1)$
    \item $(1,2), (-1,-2)$
    \item $(1,2), (2,1)$
    \item $(1,4), (4,1)$
    \item $(1,2), (2,1)$
    \item $(3,-1), (-1,3)$
\end{enumerate}

\end{document}